% A.127: all finite harmonic exponents divisible by four.
\subsection{Harmonic initialization at every finite Kummer exponent}

\begin{statement}{A.127}{Proved}
{All harmonic common characters vanish at every exponent divisible by four}
Let $4\mid N$, put $\mathcal C=\{0,1,\tfrac12\}$,
and take $R\in\Q\setminus\mathcal C$.
Consider the actual pointed cover and boundary fibre
\[
 \begin{gathered}
 C_R^{(N)}:\quad T-a=(R-a)X_a^N\quad(a\in\mathcal C),\\
 \mathcal E_{N,c}
   =\Q[Z_a^{\pm1}:a\in\mathcal C]
       \big/\left(Z_a^N-\frac{c-a}{R-a}\right)_{a\in\mathcal C},
 \qquad c\in\Q\setminus(\mathcal C\cup\{R\}).
 \end{gathered}
\tag{A.127.1}
\]
Let $B_N$ be the full common abelian field of this
fibre, as computed by A.126.
For every character $\chi$ of $\operatorname{Gal}(B_N/\Q)$,
every place $v$, and every ordinary point of
$C_R^{(N)}(\Q_v)$ with $T\ne c$, one has
\[
                  \left(\frac{T-c}{R-c},\chi\right)=0
                                      \quad\text{in }\operatorname{Br}(\Q_v).
\tag{A.127.2}
\]
Consequently, for any rational new puncture
$c\notin\mathcal C$, any finite set $S$ of places
and one prescribed ordinary point of $C_R^{(N)}$
at each $v\in S$ avoiding $T=c$, there is a rational
$a_c\ne0$ for which
\[
 \begin{gathered}
 T-a=(R-a)X_a^N\quad(a\in\mathcal C),\qquad
 T-c=a_cY^N,\\
          (T,(X_a)_a,Y)\longmapsto(T,(X_a)_a)
 \end{gathered}
\tag{A.127.3}
\]
has ordinary points over every completion and
lifts all the prescribed points.
This gives finite initialization at every exponent
in the fully cofinal set $\{N:4\mid N\}$.
It does not select these new coefficients compatibly
under divisibility of $N$ or retain an infinite
prescribed tuple of old local points.
\end{statement}
\begin{proof}
\emph{Reduction through the actual odd-factor map.}
Write $N=nv$, where $n=2^r\ge4$ and $v$ is odd.
The map
\[
 C_R^{(N)}\longrightarrow C_R^{(n)},\qquad
                       X_a\longmapsto X_a^v
\tag{A.127.4}
\]
retains $T,R,c$ and the normalized function
$f=(T-c)/(R-c)$.
A.126.21 identifies the common fields for the same
radicands exactly: $B_N=B_n$.
Its proof uses the actual upper root stabilizers,
not an equality of field degrees.
Therefore (A.127.2) at $N$ is precisely its value
at the image (A.127.4), with the same character.
We prove it for every dyadic $n$.

Put $D_a=(c-a)/(R-a)$.
A.57, or the ambient field in A.126, gives
\[
B_n\subset\Q(\mu_{2n},\sqrt P:P\in\mathcal P),
 \qquad
              \operatorname{exp}\operatorname{Gal}(B_n/\Q)\mid n/2.
\tag{A.127.5}
\]
Here $\mathcal P$ consists of the primes $P$ with
$v_P(D_a)\ne0$ for at least one $a$.
Indeed the cyclotomic group has exponent $n/2$
and the additional quadratic factors have exponent
two, which divides $n/2$.
Thus every local character field has degree
dividing $n/2$.

The actual harmonic root maps in (A.123.7) hold
with exponent $n$ without any change.
On the boundary fibres they are the same
invertible monomial maps
\[
 (Z_0,Z_1,Z_{1/2})\longmapsto
 (Z_1,Z_0,Z_{1/2}),\qquad
 (Z_0^{-1},Z_{1/2}/Z_0,Z_1/Z_0).
\tag{A.127.6}
\]
The first map leaves $f$ unchanged; the second
multiplies it by $R/T=X_0^{-n}$.
They identify every field factor and the actual
common characters.
The multiplier is a norm from each character
field by (A.127.5).
These maps therefore preserve the evaluation.

\emph{The exact dyadic power and cyclotomic norm groups.}
For $j\ge2$ the squaring map takes
$1+2^j\Z_2$ onto $1+2^{j+1}\Z_2$.
The normalized polynomial is
\[
 (1+2^jz)^2
       =1+2^{j+1}\bigl(z+2^{j-1}z^2\bigr);
\tag{A.127.7}
\]
its derivative is odd, so each target has a unique
root on this branch by Hensel lifting.
Every odd unit is a sign times an element of
$1+4\Z_2$.
Iterating (A.127.7) proves, for every $b\ge1$,
\[
                    (\Z_2^\times)^{2^b}
                              =1+2^{b+2}\Z_2.
\tag{A.127.8}
\]

We also retain the full signed cyclotomic norm map.
For $r\ge2$, the local reciprocity map on
$\Q_2(\mu_{2^r})$ is
\[
       u\longmapsto(\zeta\longmapsto\zeta^{u^{-1}})
                                      \quad(u\in\Z_2^\times).
\tag{A.127.9}
\]
This is the explicit unit component of
\cite[Theorem 4.1.2 and Example 4.1.4]{BIR_KedlayaCFT}.
Its norm kernel on units is exactly
$1+2^r\Z_2$.
In particular the sign is retained: $-1$ acts
by inversion and is excluded from that kernel.
If all characters of this cyclotomic field have
zero cyclic evaluation on a unit $f$, their
perfect pairing with its reciprocity image makes
that image the identity.
Thus
\[
 \text{all cyclotomic evaluations on }f\text{ vanish}
        \quad\Longrightarrow\quad f\in1+2^r\Z_2.
\tag{A.127.10}
\]

\emph{Induction on the dyadic exponent.}
At $n=4$, A.112 proves the assertion for every
common character.
Let $n=2^r$ with $r\ge3$, and assume the assertion
at $n/2$.
The coordinate-square map
$C_R^{(n)}\to C_R^{(n/2)}$ retains $f$, so every
character of $B_{n/2}$ has zero evaluation at our
point.
If $\mu_n\not\subset B_{n/2}$, the exact stopping
rule A.126 gives $B_n=B_{n/2}$ and finishes this
step.
We may assume
\[
                       \mu_n\subset B_{n/2}.
\tag{A.127.11}
\]
Then $i\in K:=\Q(\sqrt{D_a}:a\in\mathcal C)$.
Indeed if $i\notin K$, A.80 would give
$B_{n/2}=K$ (or the same base equality at $n/2=2$),
contradicting (A.127.11).
The individual-squareclass pairing therefore
supplies a subset whose product of the $D_a$
is a negative rational square.

Use (A.127.6) to arrange
$e=v_2(R)\ge1$, exactly as in A.123.
Write $e_T=v_2(T)$ and $s=v_2(T-R)$.
The old $n$th-power ratios give
\[
 \begin{gathered}
 v_2(T)\equiv e,\quad v_2(T-1)\equiv0,\quad
                 v_2(T-\tfrac12)\equiv-1\pmod n,\\
 e_T>0,\qquad e_T\equiv e\pmod n,\qquad s\ge r+2.
 \end{gathered}
\tag{A.127.12}
\]
For negative $e_T$, one has $v_2(T-1)=e_T$,
so the second congruence first gives $n\mid e_T$.
Then $e_T\le-n<-1$, and $v_2(T-\tfrac12)=e_T$
contradicts the third congruence.
For $e_T=0$, both $e$ and $v_2(T-1)$ would be
positive multiples of $n$.
The ratio $(1-2T)/(1-2R)$ would then be $-1$
modulo $2^{r+2}$.
It must instead be $1$ modulo $2^{r+2}$ by
(A.127.8).
Here $n\ge r+2$, so the asserted congruence
follows from those depths.
Finally $(1-T)/(1-R)$ is an $n$th-power unit;
(A.127.8) gives the bound on $s$.

If $e_T=e$, all old valuation coordinates agree.
With $M=\max_a v_2(R-a)=e$, the same unit equations
give $s\ge M+r+2$.
The negative-square subset just obtained forces
$v_2(c-R)\le M+1$:
otherwise every $D_a$ lies in $1+4\Z_2$ and
no subset product can be a negative unit square.
Consequently
\[
                        f\in1+2^{r+1}\Z_2
                                 =(\Z_2^\times)^{n/2}.
\tag{A.127.13}
\]
It is a norm from every character field in
(A.127.5), proving this case.
This argument also proves the same equal-valuation
calculation for any finite set of old punctures,
with its displayed maximum $M$.

For unequal depths, the exact consequences of
(A.127.12) are
\[
 e,e_T\ge r+2,\qquad
 s=\min(e,e_T),\qquad e\equiv e_T\pmod n.
\tag{A.127.14}
\]
Put $k=v_2(c)$.
When $k\le0$, $v_2(R-c)=k$ and
$v_2(f-1)=s-k\ge r+2$, so (A.127.13) again applies.

Suppose $1\le k\le r+1$.
Then $f$ is a unit and
\[
                       v_2(f-1)=s-k.
\tag{A.127.15}
\]
Every cyclotomic character of $\Q(\mu_n)$ belongs
to $B_{n/2}$ by (A.127.11).
The inductive vanishing and (A.127.10) force
$f\in1+2^r\Z_2$.
If the depth in (A.127.15) is at least $r+1$,
$f$ is an $(n/2)$th power and the proof is finished.

The only remaining depth is therefore
\[
                s-k=r,\qquad k\ge2,\qquad
                        e-k\equiv r\pmod n.
\tag{A.127.16}
\]
The function $f$ is an actual $(n/4)$th power
by (A.127.8).
We now use every relation in the complete
A.126 formula.
For its representatives $0\le t_a<n/2$, one has
\[
 \prod_aD_a^{t_a}=(-1)^{\delta_t}r_t^{\,n/2},
             \qquad r_t\in\Q_{>0}.
\tag{A.127.17}
\]
Only $D_0$ has nonzero dyadic valuation; the
others satisfy
\[
 v_2(D_0)=k-e,\qquad
 D_1,D_{1/2}\in1+4\Z_2.
\tag{A.127.18}
\]
Taking valuations in (A.127.17) gives
\[
              2^{r-1}\mid t_0(k-e).
\tag{A.127.19}
\]
Since $k-e\equiv-r\pmod{2^r}$, one has
$v_2(k-e)=v_2(r)<r-1$ for every $r\ge3$.
Thus $t_0$ is even.

Write $D_0=2^{k-e}u$ with $u$ odd, and remove
the equal powers of two from both sides of
(A.127.17).
Its normalized left unit is
$u^{t_0}D_1^{t_1}D_{1/2}^{t_{1/2}}$, which is
$1$ modulo four because $t_0$ is even and
(A.127.18) holds.
A negative right side would be $-1$ modulo
four: $n/2$ is even and the normalized unit
of $r_t$ is odd.
Therefore every relation in (A.127.17) has
$\delta_t=0$.

The actual formula A.126 now says that $B_n$
is generated by $B_{n/2}$ and quadratic fields.
All characters of $B_{n/2}$ have order dividing
$n/4$, by (A.127.5) at the preceding exponent.
Since $n/4\ge2$, the added quadratic characters
have this same exponent bound.
Thus every character of $B_n$ has order dividing
$n/4$, and the actual $(n/4)$th power $f$ kills
every evaluation.
This proves the remaining shallow-$c$ case
without omitting any mixed relation.

Finally suppose $e,k\ge r+2$.
The other two radicands
\[
              D_1=\frac{1-c}{1-R},\qquad
              D_{1/2}=\frac{1-2c}{1-2R}
\tag{A.127.20}
\]
lie in $1+2^{r+2}\Z_2$ and hence are actual
$n$th powers.
Choose their roots in $\Q_2$.
Specializing those coordinates gives the exact
surjection
\[
     \mathcal E_{n,c}\otimes\Q_2
                  \longrightarrow\Q_2[Z_0]/(Z_0^n-D_0).
\tag{A.127.21}
\]
Every target factor contains the local cyclic
field of each common character.
For one root $\theta_E$ in each factor $E$,
\[
       f=\frac{X_0^n-D_0}{1-D_0}
         =\prod_E N_{E/\Q_2}
                    \left(\frac{X_0-\theta_E}{1-\theta_E}\right).
\tag{A.127.22}
\]
These are actual polynomial norms and all
denominators are nonzero.
Norm transitivity proves every desired
evaluation.
Equations (A.127.13)--(A.127.22) complete the
dyadic induction for all $r\ge2$.

\emph{Odd places with the same field-factor maps.}
Let $p$ be odd.
The four punctures have distinct reductions and
$n$ is invertible in $\Z_p$.
Every principal unit is an $n$th power, by
Hensel applied to $Z^n-u$ at its residue root.
The maps (A.127.6) preserve and permute the
boundary discs exactly as in A.123.
The following list retains every configuration.

If $R,c$ share a boundary disc, send it to
zero.
Then $D_1,D_{1/2}$ are principal units, and
specialization to their actual $n$th roots
gives (A.127.21)--(A.127.22) over $\Q_p$.
If $R,T$ share a boundary disc and $c$ lies
elsewhere, the same normalization gives
$f=(1-T/c)/(1-R/c)$, a principal unit and
therefore an $n$th power.

When exactly one of $R,T$ is good, exchange
them by dividing each boundary root by its
actual old root; this changes the normalized
entry to $f^{-1}$ and leaves the local common
character unchanged.
Thus take $R$ good.
For $c$ in a boundary disc, send that disc to
zero.
If $T$ shares it, the first case applies after
an exchange.
Otherwise
\[
                     f/(T/R)=\frac{1-c/T}{1-c/R}
                                        \in1+p\Z_p,
\tag{A.127.23}
\]
so $f$ is an $n$th power.
For $R,c$ both good, all $D_a$ are units, so
their $n$th roots lie in an unramified extension:
lift their residue roots over a sufficiently
large finite residue field.
If $R$ or a good $T$ reduces to $c$, each $D_a$
is an actual $n$th power, and the fibre has
a local rational root tuple; every common
character is then locally zero.
In every other case $v_p(f)=0$, except for
$T$ in the disc of infinity, where
$(T-1)/(R-1)$ gives $v_p(f)=v_p(T)\in n\Z$.
The local character field is unramified with
degree dividing $n/2$, so $f$ is a norm.
This norm assertion follows from the finite-field
norm, principal-unit roots and the norm of
the base uniformizer, as in A.123.

It remains that $R,T$ lie in different boundary
discs.
A $c$ in either disc is the first case.
For $c$ in a third boundary disc, send it to
zero and apply (A.127.23).
For good $c$, send $R$ to zero.
If $T$ lies in another finite boundary disc,
$T/R$ forces $v_p(R)\in n\Z$.
If $T$ is in the disc of infinity, the ratio
$(T-1)/(R-1)$ first forces $v_p(T)\in n\Z$
and then $T/R$ forces $v_p(R)\in n\Z$.
All the radicands consequently have valuations
in $n\Z$ and have roots in an unramified
extension after scaling by integral uniformizer
powers.
Also $v_p(f)$ is zero or $v_p(T)\in n\Z$.
The same unramified norm map proves vanishing.
This exhausts the odd places.

\emph{The real place and every finite exponent.}
The positive old ratios place $R,T$ in the same
interval between the real old punctures.
If $c$ is outside their intervening interval,
then $f>0$, which is a norm from $\C$.
If $c$ lies between them, every $D_a>0$ and
the boundary has a real $n$th-root tuple;
its common characters are locally trivial.
This proves the dyadic-exponent assertion at
all places.
Equation (A.127.4) and the exact odd-factor
field equality prove it for every $4\mid N$.

Finally apply the precise finite criterion A.85.
For $c\ne R$, subtracting the global constant
$(R-c,\chi)$ leaves the locally zero classes
(A.127.2), while that constant's invariants
sum to zero.
The criterion supplies the rational coefficient
and the actual map (A.127.3), retaining the
finite prescribed collection of old points.
For $c=R$, the rational boundary root tuple
$(1,1,1)$ makes the common field $\Q$ and the
criterion has no nonzero characters.
This proves exactly the stated finite
initialization for every exponent divisible
by four.
\end{proof}

The retained certificate
\path{code/harmonic_all_exponents_certificate.py}
checks the normalized squaring maps, finite
cyclotomic norm kernels including their sign,
all tested depth and mixed-relation implications,
the polynomial norm maps, and literal odd-factor
CRT root lifts.
Its output accompanies the code.
The induction and its field-factor comparisons
prove the unrestricted exponent assertion.

